\documentclass[12pt,a4paper]{article}
\usepackage[T1]{fontenc}
\usepackage[utf8]{inputenc}
\usepackage[magyar]{babel}
\usepackage{amsmath}
\usepackage{amssymb}
\usepackage{geometry}
\geometry{margin=2.5cm}

\title{Ellenfél alapú támadások (FGSM, PGD) implementációja\\
konvolúciós neurális hálózatokhoz\\
interaktív tesztkörnyezetben}
\author{Gabrity Gábor\\Neptun kód: D9V09Z}
\date{\today}

\begin{document}

\maketitle
\newpage

\section*{Nyilatkozat}

Alulírott Gabrity Gábor (Neptun kód: D9V09Z), az Eötvös Loránd Tudományegyetem Informatikai Karának
hallgatója kijelentem, hogy ezt a szakdolgozatot saját magam készítettem, és abban csak a megadott
forrásokat használtam fel. Minden olyan részt, amelyet szó szerint vagy azonos tartalommal, de
átfogalmazva más forrásból átvettem, egyértelműen, a forrás megadásával megjelöltem.

A dolgozat bírálója által hátrányos értékeléssel nem járó exkluzív elektronikus hozzáférését a
benyújtástól számított két évig nem kérem.

\vspace{2cm}

Budapest, \today

\vspace{2cm}
\begin{flushright}
Gabrity Gábor s.k.\\
hallgató aláírása
\end{flushright}

\newpage

\begin{abstract}
A mesterséges intelligencia rohamos fejlődésével párhuzamosan kritikus kérdéssé vált a modellek biztonsága
és megbízhatósága. Kutatások kimutatták, hogy a modern konvolúciós neurális hálózati architektúrák
sebezhetőek az ellenfél alapú támadásokkal szemben, mint például az FGSM (Fast Gradient Sign Method) és a
PGD (Projected Gradient Descent). Ezek az algoritmusok a bemeneti adatot képesek olyan módon -- az emberi szem számára
szinte észrevétlenül -- módosítani, hogy a nagy megbízhatóságú modelleket is megtévesztik.

Szakdolgozatom célja egy olyan interaktív tesztkörnyezet fejlesztése, amely grafikus felhasználói felületen keresztül
teszi lehetővé ezen biztonsági rések vizuális elemzését. A projekt fókuszában az AI Red Teaming módszertan áll: a
fejlesztő vagy biztonsági szakember a szoftver segítségével szisztematikusan ,,támadja'' a választott
képfelismerő modellt, hogy feltárja annak gyenge pontjait.

A keretrendszer több elismert támadási algoritmust implementál, köztük a gyors, egylépéses FGSM-et és a robusztusabb,
iteratív PGD-t. A felhasználó számára lehetőség nyílik mind általános tévesztés (non-targeted), mind specifikus
célosztályra irányuló (targeted) támadások végrehajtására, miközben interaktív módon, csúszkák segítségével
szabályozhatja a perturbáció mértékét és az iterációk számát.

Az alkalmazás hármas vizuális megjelenítést alkalmaz: megjeleníti az eredeti képet, a generált zajt (perturbációt)
és a végső, támadott képet, miközben mutatja a modell predikcióinak változását. A rendszer kiegészülhet egy
modellvédelmi modullal is, amely különböző képfeldolgozási szűrők -- például homályosítás -- segítségével
próbálja megsemmisíteni a beágyazott zajt, mérve ezzel a modell védekezőképességét.
\end{abstract}
\newpage

\tableofcontents
\newpage

\section{Ellenfél alapú támadások}

A mély neurális hálózatok -- különösen a képfelismerésben alkalmazott konvolúciós architektúrák -- rendkívül
érzékenyek a bemeneti adatok apró, célzott módosításaira. Az ilyen módosításokat \textit{ellenfél alapú
perturbációknak} (adversarial perturbations) nevezzük, az ezeket előállító eljárásokat pedig \textit{ellenfél
alapú támadásoknak} (adversarial attacks)~\cite{goodfellow2015explaining}.

A támadások alapelve a következő: adott egy $f$ osztályozó modell, egy $x$ eredeti bemenet és annak helyes
$y$ címkéje. A cél egy olyan $x_{\mathrm{adv}}$ módosított bemenet előállítása, amelyre
$f(x_{\mathrm{adv}}) \neq y$, miközben a perturbáció $\delta = x_{\mathrm{adv}} - x$ valamilyen norma
szerint kicsi marad: $\|\delta\|_p \leq \epsilon$. A gyakorlatban leggyakrabban az $\ell_\infty$ normát
alkalmazzák, amely a pixelenkénti maximális eltérést korlátozza.

A következő alfejezetekben a két legelterjedtebb gradiens alapú támadási módszert mutatom be részletesen.

\subsection{FGSM (Fast Gradient Sign Method)}

Az FGSM algoritmust Goodfellow és társai javasolták 2015-ben~\cite{goodfellow2015explaining}. Az eljárás
egyetlen lépésben állítja elő a perturbációt a veszteségfüggvény bemenetre vonatkozó gradiensének
előjelét felhasználva.

\subsubsection{Matematikai háttér}

Legyen $J(\theta, x, y)$ a modell veszteségfüggvénye (például keresztentrópia), ahol $\theta$ a hálózat
tanult paraméterei, $x$ a bemeneti kép és $y$ az eredeti címke. Az FGSM a veszteségfüggvény
$x$-re vonatkozó gradiensét számítja ki visszaterjesztéssel (backpropagation), majd a perturbációt
a gradiens előjelének $\epsilon$-szorosával képezi:
\[
  x_{\mathrm{adv}} = x + \epsilon \cdot \mathrm{sign}\bigl(\nabla_x J(\theta, x, y)\bigr)
\]

A $\mathrm{sign}(\cdot)$ függvény elemenként az előjelet adja vissza ($-1$, $0$ vagy $+1$), így a
perturbáció minden pixel esetén pontosan $\pm\epsilon$ lesz. Ez biztosítja, hogy a módosítás az
$\ell_\infty$ norma szerint pontosan $\epsilon$ nagyságú.

\subsubsection{Az algoritmus lépései}

\begin{enumerate}
  \item Előreterjesztés (forward pass): a modell kiszámítja az $x$ bemenethez tartozó predikciót és a
        veszteségfüggvény $J(\theta, x, y)$ értékét.
  \item Visszaterjesztés (backward pass): kiszámítjuk a gradienst $\nabla_x J(\theta, x, y)$.
  \item Perturbáció előállítása: $\delta = \epsilon \cdot \mathrm{sign}\bigl(\nabla_x J(\theta, x, y)\bigr)$.
  \item Támadott kép: $x_{\mathrm{adv}} = \mathrm{clip}(x + \delta,\; 0,\; 1)$, ahol a $\mathrm{clip}$
        függvény biztosítja, hogy a pixelértékek a $[0, 1]$ tartományban maradjanak.
\end{enumerate}

\subsubsection{Célzott változat}

Az alapváltozat (\textit{non-targeted}) a helyes osztály veszteségét maximalizálja. A \textit{targeted}
változatban egy $y_{\mathrm{target}}$ célosztályt adunk meg, és a perturbációt úgy állítjuk elő, hogy a
modell ebbe az osztályba sorolja a képet:
\[
  x_{\mathrm{adv}} = x - \epsilon \cdot \mathrm{sign}\bigl(\nabla_x J(\theta, x, y_{\mathrm{target}})\bigr)
\]
Az előjel megfordul, mert ebben az esetben a célosztályhoz tartozó veszteséget \textit{minimalizáljuk}.

\subsubsection{Előnyök és korlátok}

Az FGSM fő előnye a sebessége: egyetlen gradiens-számítás elegendő. Hátránya, hogy egyetlen nagy lépést tesz,
így kisebb $\epsilon$ értékek mellett gyakran kevésbé hatékony, mint az iteratív módszerek. Emellett a
perturbáció iránya nem feltétlenül optimális, csupán a gradiens előjelét követi.

\subsection{PGD (Projected Gradient Descent)}

A PGD algoritmust Madry és társai vezették be 2018-ban~\cite{madry2018towards} az FGSM iteratív
általánosításaként. A módszert a szakirodalom gyakran a ,,legerősebb elsőrendű támadásnak'' tekinti
az $\ell_\infty$ korlátos perturbációk között.

\subsubsection{Matematikai háttér}

A PGD a következő iteratív frissítési szabályt alkalmazza, $t = 0, 1, \ldots, T-1$ lépésben:
\[
  x^{(t+1)} = \Pi_{\mathcal{B}_\epsilon(x)} \Bigl(
    x^{(t)} + \alpha \cdot \mathrm{sign}\bigl(\nabla_x J(\theta, x^{(t)}, y)\bigr)
  \Bigr)
\]
ahol:
\begin{itemize}
  \item $\alpha$ a lépésnagyság (step size), jellemzően $\alpha < \epsilon$,
  \item $T$ az iterációk száma,
  \item $\Pi_{\mathcal{B}_\epsilon(x)}$ a projekciós operátor, amely a módosított mintát visszavetíti
        az eredeti $x$ körüli $\epsilon$ sugarú $\ell_\infty$ gömbbe:
        \[
          \Pi_{\mathcal{B}_\epsilon(x)}(z) = \mathrm{clip}(z,\; x - \epsilon,\; x + \epsilon)
        \]
\end{itemize}

\subsubsection{Az algoritmus lépései}

\begin{enumerate}
  \item Kezdeti pont megválasztása: $x^{(0)}$ lehet maga az eredeti $x$, vagy egy véletlen pont az
        $\mathcal{B}_\epsilon(x)$ környezetből (random restart).
  \item Iteráció ($t = 0, \ldots, T-1$):
    \begin{enumerate}
      \item Előreterjesztés és veszteségszámítás az aktuális $x^{(t)}$ pontra.
      \item Gradiens számítása: $g^{(t)} = \nabla_x J(\theta, x^{(t)}, y)$.
      \item Lépés a gradiens irányába: $\hat{x}^{(t+1)} = x^{(t)} + \alpha \cdot \mathrm{sign}(g^{(t)})$.
      \item Projekció: $x^{(t+1)} = \Pi_{\mathcal{B}_\epsilon(x)}(\hat{x}^{(t+1)})$.
      \item Pixelértékek vágása a $[0, 1]$ tartományba.
    \end{enumerate}
  \item A végső támadott kép: $x_{\mathrm{adv}} = x^{(T)}$.
\end{enumerate}

\subsubsection{Célzott változat}

A targeted PGD a célzott FGSM-hez hasonlóan a célosztály veszteségét minimalizálja:
\[
  x^{(t+1)} = \Pi_{\mathcal{B}_\epsilon(x)} \Bigl(
    x^{(t)} - \alpha \cdot \mathrm{sign}\bigl(\nabla_x J(\theta, x^{(t)}, y_{\mathrm{target}})\bigr)
  \Bigr)
\]

\subsubsection{Hiperparaméterek hatása}

A PGD viselkedését három fő hiperparaméter határozza meg:
\begin{itemize}
  \item $\epsilon$: a megengedett maximális perturbáció. Nagyobb $\epsilon$ erősebb támadást tesz lehetővé,
        de a módosítás emberi szemmel is láthatóvá válhat.
  \item $\alpha$: a lépésnagyság. Túl kis $\alpha$ lassú konvergenciát eredményez, túl nagy $\alpha$ pedig
        instabilitáshoz vezethet. Tipikus választás: $\alpha = \epsilon / T$ vagy $\alpha = 2.5 \cdot \epsilon / T$.
  \item $T$: az iterációk száma. Több iteráció általában erősebb támadást eredményez, de növeli a
        számítási költséget.
\end{itemize}

\subsubsection{Előnyök és korlátok}

A PGD az FGSM-nél lényegesen erősebb támadásokat képes generálni, mivel több kis lépésben közelíti
a lokális optimumot. A random restart technikával kombinálva különösen hatékony. Hátránya, hogy
$T$-szer annyi gradiens-számítást igényel, mint az FGSM, így nagyobb a számítási költsége.
\newpage

\begin{thebibliography}{9}
\bibitem{goodfellow2015explaining}
Ian J. Goodfellow, Jonathon Shlens, Christian Szegedy.
\textit{Explaining and Harnessing Adversarial Examples}.
arXiv:1412.6572, 2015.

\bibitem{madry2018towards}
Aleksander Madry, Aleksandar Makelov, Ludwig Schmidt, Dimitris Tsipras, Adrian Vladu.
\textit{Towards Deep Learning Models Resistant to Adversarial Attacks}.
arXiv:1706.06083, 2018.
\end{thebibliography}

\end{document}
